\hypertarget{dictionnaires}{%
\section{Dictionnaires}\label{dictionnaires}}

Un \textbf{dictionnaire} est un type de données intégré à Python qui
ressemble à une liste, mais il est plus général. Dans une liste, les
indices doivent être des entiers ; dans un dictionnaire, les indices
peuvent être de presque n\textquotesingle importe quel type.

Un dictionnaire contient une collection de \textbf{indices}, appelés
\textbf{clés}, et chaque clé est associée à une \textbf{valeur}.
L\textquotesingle association d\textquotesingle une clé et
d\textquotesingle une valeur est appelée une \textbf{paire clé-valeur}
ou un élément.

La fonction \texttt{dict} crée un nouveau dictionnaire sans aucun
élément. Comme les crochets pour les listes, les accolades (\texttt{\{}
et \texttt{\}}) représentent un dictionnaire littéral.

\begin{Shaded}
\begin{Highlighting}[]
\OperatorTok{\textgreater{}\textgreater{}\textgreater{}}\NormalTok{ anglais\_francais }\OperatorTok{=} \BuiltInTok{dict}\NormalTok{()}
\OperatorTok{\textgreater{}\textgreater{}\textgreater{}}\NormalTok{ anglais\_francais}
\NormalTok{\{\}}
\OperatorTok{\textgreater{}\textgreater{}\textgreater{}}\NormalTok{ eng2fr }\OperatorTok{=}\NormalTok{ \{}\StringTok{\textquotesingle{}un\textquotesingle{}}\NormalTok{: }\StringTok{\textquotesingle{}one\textquotesingle{}}\NormalTok{, }\StringTok{\textquotesingle{}deux\textquotesingle{}}\NormalTok{: }\StringTok{\textquotesingle{}two\textquotesingle{}}\NormalTok{, }\StringTok{\textquotesingle{}trois\textquotesingle{}}\NormalTok{: }\StringTok{\textquotesingle{}three\textquotesingle{}}\NormalTok{\}}
\OperatorTok{\textgreater{}\textgreater{}\textgreater{}}\NormalTok{ eng2fr}
\NormalTok{\{}\StringTok{\textquotesingle{}un\textquotesingle{}}\NormalTok{: }\StringTok{\textquotesingle{}one\textquotesingle{}}\NormalTok{, }\StringTok{\textquotesingle{}deux\textquotesingle{}}\NormalTok{: }\StringTok{\textquotesingle{}two\textquotesingle{}}\NormalTok{, }\StringTok{\textquotesingle{}trois\textquotesingle{}}\NormalTok{: }\StringTok{\textquotesingle{}three\textquotesingle{}}\NormalTok{\}}
\end{Highlighting}
\end{Shaded}

L\textquotesingle ordre des paires clé-valeur n\textquotesingle est pas
prévisible. Pour afficher un élément, vous utilisez une clé entre
crochets :

\begin{Shaded}
\begin{Highlighting}[]
\OperatorTok{\textgreater{}\textgreater{}\textgreater{}} \BuiltInTok{print}\NormalTok{(eng2fr[}\StringTok{\textquotesingle{}deux\textquotesingle{}}\NormalTok{])}
\CommentTok{\textquotesingle{}two\textquotesingle{}}
\end{Highlighting}
\end{Shaded}

Si la clé n\textquotesingle est pas dans le dictionnaire, vous obtenez
une exception :

\begin{Shaded}
\begin{Highlighting}[]
\OperatorTok{\textgreater{}\textgreater{}\textgreater{}} \BuiltInTok{print}\NormalTok{(eng2fr[}\StringTok{\textquotesingle{}quatre\textquotesingle{}}\NormalTok{])}
\PreprocessorTok{KeyError}\NormalTok{: }\StringTok{\textquotesingle{}quatre\textquotesingle{}}
\end{Highlighting}
\end{Shaded}

La fonction \texttt{len} renvoie le nombre de paires clé-valeur :

\begin{Shaded}
\begin{Highlighting}[]
\OperatorTok{\textgreater{}\textgreater{}\textgreater{}} \BuiltInTok{len}\NormalTok{(eng2fr)}
\DecValTok{3}
\end{Highlighting}
\end{Shaded}

L\textquotesingle opérateur \texttt{in} vous indique si une clé apparaît
dans le dictionnaire :

\begin{Shaded}
\begin{Highlighting}[]
\OperatorTok{\textgreater{}\textgreater{}\textgreater{}} \StringTok{\textquotesingle{}un\textquotesingle{}} \KeywordTok{in}\NormalTok{ eng2fr}
\VariableTok{True}
\OperatorTok{\textgreater{}\textgreater{}\textgreater{}} \StringTok{\textquotesingle{}one\textquotesingle{}} \KeywordTok{in}\NormalTok{ eng2fr}
\VariableTok{False}
\end{Highlighting}
\end{Shaded}

L\textquotesingle opérateur \texttt{in} recherche parmi les clés, pas
les valeurs. Pour vérifier si une valeur apparaît dans un dictionnaire,
vous pouvez utiliser la méthode \texttt{values}, qui renvoie les valeurs
sous forme de collection, puis utiliser l\textquotesingle opérateur
\texttt{in} :

\begin{Shaded}
\begin{Highlighting}[]
\OperatorTok{\textgreater{}\textgreater{}\textgreater{}}\NormalTok{ valeurs }\OperatorTok{=}\NormalTok{ eng2fr.values()}
\OperatorTok{\textgreater{}\textgreater{}\textgreater{}} \StringTok{\textquotesingle{}one\textquotesingle{}} \KeywordTok{in}\NormalTok{ valeurs}
\VariableTok{True}
\end{Highlighting}
\end{Shaded}

Le dictionnaire comme ensemble de compteurs
-\/-\/-\/-\/-\/-\/-\/-\/-\/-\/-\/-\/-\/-\/-\/-\/-\/-\/-\/-\/-\/-\/-\/-\/-\/-\/-\/-\/-\/-\/-\/-\/-\/-\/-\/-\/-\/-\/-\/-\/-\/-

Supposons que vous receviez une chaîne de caractères et que vous
souhaitiez compter combien de fois chaque lettre y apparaît. Une façon
de faire est de créer un dictionnaire avec les lettres comme clés et les
compteurs comme valeurs.

\begin{Shaded}
\begin{Highlighting}[]
\KeywordTok{def}\NormalTok{ histogramme(s):}
\NormalTok{    d }\OperatorTok{=} \BuiltInTok{dict}\NormalTok{()}
    \ControlFlowTok{for}\NormalTok{ c }\KeywordTok{in}\NormalTok{ s:}
        \ControlFlowTok{if}\NormalTok{ c }\KeywordTok{not} \KeywordTok{in}\NormalTok{ d:}
\NormalTok{            d[c] }\OperatorTok{=} \DecValTok{1}
        \ControlFlowTok{else}\NormalTok{:}
\NormalTok{            d[c] }\OperatorTok{+=} \DecValTok{1}
    \ControlFlowTok{return}\NormalTok{ d}
\end{Highlighting}
\end{Shaded}

La fonction \texttt{histogramme} prend une chaîne de caractères et
renvoie un dictionnaire contenant les caractères uniques comme clés et
leurs fréquences comme valeurs.

Les dictionnaires possèdent une méthode appelée \texttt{get} qui prend
une clé et une valeur par défaut. Si la clé est dans le dictionnaire,
\texttt{get} renvoie la valeur correspondante ; sinon, elle renvoie la
valeur par défaut. Par exemple :

\begin{Shaded}
\begin{Highlighting}[]
\KeywordTok{def}\NormalTok{ histogramme(s):}
\NormalTok{    d }\OperatorTok{=} \BuiltInTok{dict}\NormalTok{()}
    \ControlFlowTok{for}\NormalTok{ c }\KeywordTok{in}\NormalTok{ s:}
\NormalTok{        d[c] }\OperatorTok{=}\NormalTok{ d.get(c, }\DecValTok{0}\NormalTok{) }\OperatorTok{+} \DecValTok{1}
    \ControlFlowTok{return}\NormalTok{ d}
\end{Highlighting}
\end{Shaded}

\hypertarget{boucles-et-dictionnaires}{%
\subsection{Boucles et dictionnaires}\label{boucles-et-dictionnaires}}

Si vous utilisez une instruction \texttt{for} avec un dictionnaire, elle
parcourt les clés du dictionnaire. Par exemple, cette fonction imprime
chaque clé et sa valeur associée :

\begin{Shaded}
\begin{Highlighting}[]
\KeywordTok{def}\NormalTok{ imprimer\_histogramme(h):}
    \ControlFlowTok{for}\NormalTok{ c }\KeywordTok{in}\NormalTok{ h:}
        \BuiltInTok{print}\NormalTok{(c, h[c])}
\end{Highlighting}
\end{Shaded}

Les clés ne sont pas dans un ordre trié. Pour parcourir les clés dans un
ordre trié, vous pouvez utiliser la fonction \texttt{sorted} :

\begin{Shaded}
\begin{Highlighting}[]
\ControlFlowTok{for}\NormalTok{ c }\KeywordTok{in} \BuiltInTok{sorted}\NormalTok{(h):}
    \BuiltInTok{print}\NormalTok{(c, h[c])}
\end{Highlighting}
\end{Shaded}

\hypertarget{recherche-inversuxe9e}{%
\subsection{Recherche inversée}\label{recherche-inversuxe9e}}

Étant donné un dictionnaire \texttt{d} et une clé \texttt{k}, il est
facile de trouver la valeur correspondante : \texttt{v\ =\ d{[}k{]}}.
C\textquotesingle est une \textbf{recherche}.

Mais que faire si vous avez une valeur et que vous voulez trouver la clé
correspondante ? C\textquotesingle est une \textbf{recherche inversée}
(\emph{reverse lookup}).

\begin{Shaded}
\begin{Highlighting}[]
\KeywordTok{def}\NormalTok{ recherche\_inverse(d, v):}
    \ControlFlowTok{for}\NormalTok{ k }\KeywordTok{in}\NormalTok{ d:}
        \ControlFlowTok{if}\NormalTok{ d[k] }\OperatorTok{==}\NormalTok{ v:}
            \ControlFlowTok{return}\NormalTok{ k}
        \ControlFlowTok{raise} \PreprocessorTok{LookupError}\NormalTok{(}\StringTok{\textquotesingle{}La valeur ne redevient pas une clé dans le dictionnaire.\textquotesingle{}}\NormalTok{)}
\end{Highlighting}
\end{Shaded}

\hypertarget{dictionnaires-et-listes}{%
\subsection{Dictionnaires et listes}\label{dictionnaires-et-listes}}

Les listes peuvent apparaître comme des valeurs dans un dictionnaire.
Par exemple, si vous voulez inverser un histogramme (obtenir un
dictionnaire qui associe des fréquences à des listes de lettres ayant
cette fréquence), vous pouvez écrire :

\begin{Shaded}
\begin{Highlighting}[]
\KeywordTok{def}\NormalTok{ inverser\_dictionnaire(d):}
\NormalTok{    inverse }\OperatorTok{=} \BuiltInTok{dict}\NormalTok{()}
    \ControlFlowTok{for}\NormalTok{ k }\KeywordTok{in}\NormalTok{ d:}
\NormalTok{        v }\OperatorTok{=}\NormalTok{ d[k]}
        \ControlFlowTok{if}\NormalTok{ v }\KeywordTok{not} \KeywordTok{in}\NormalTok{ inverse:}
\NormalTok{            inverse[v] }\OperatorTok{=}\NormalTok{ [k]}
        \ControlFlowTok{else}\NormalTok{:}
\NormalTok{            inverse[v].append(k)}
    \ControlFlowTok{return}\NormalTok{ inverse}
\end{Highlighting}
\end{Shaded}

Les listes peuvent être des valeurs dans un dictionnaire, mais
\textbf{les listes ne peuvent pas être des clés}. Les clés doivent être
\textbf{hachables} (\emph{hashable}), ce qui signifie
qu\textquotesingle elles possèdent une fonction de hachage immuable. Les
types mutables comme les listes ne sont pas hachables.

\hypertarget{muxe9mouxefsation}{%
\subsection{Mémoïsation}\label{muxe9mouxefsation}}

Si vous avez exécuté la fonction de Fibonacci récursive, vous avez
peut-être remarqué que plus le nombre grandit, plus le temps de calcul
explose. Cela est dû au fait que la fonction recalcule les mêmes valeurs
encore et encore.

Une solution consiste à stocker les résultats déjà calculés dans un
dictionnaire, souvent appelé \textbf{mémo} (\emph{memo}). Un
dictionnaire utilisé de cette façon est un exemple de
\textbf{mémoïsation} (\emph{memoization}) :

\begin{Shaded}
\begin{Highlighting}[]
\NormalTok{connnu }\OperatorTok{=}\NormalTok{ \{}\DecValTok{0}\NormalTok{: }\DecValTok{0}\NormalTok{, }\DecValTok{1}\NormalTok{: }\DecValTok{1}\NormalTok{\}}

\KeywordTok{def}\NormalTok{ fibonacci(n):}
    \ControlFlowTok{if}\NormalTok{ n }\KeywordTok{in}\NormalTok{ connu:}
        \ControlFlowTok{return}\NormalTok{ connu[n]}
\NormalTok{    res }\OperatorTok{=}\NormalTok{ fibonacci(n}\OperatorTok{{-}}\DecValTok{1}\NormalTok{) }\OperatorTok{+}\NormalTok{ fibonacci(n}\OperatorTok{{-}}\DecValTok{2}\NormalTok{)}
\NormalTok{    connnu[n] }\OperatorTok{=}\NormalTok{ res}
    \ControlFlowTok{return}\NormalTok{ res}
\end{Highlighting}
\end{Shaded}

Variables globales -\/-\/-\/-\/-\/-\/-\/-\/-\/-\/-\/-\/-\/-\/-\/-

Dans l\textquotesingle exemple ci-dessus, la variable \texttt{connu} est
créée en dehors de la fonction, ce qui en fait une \textbf{variable
globale} (\emph{global variable}). Vous pouvez lire les variables
globales à l\textquotesingle intérieur d\textquotesingle une fonction,
mais si vous voulez les modifier ou leur réassigner une valeur, vous
devez utiliser le mot-clé \texttt{global} :

\begin{Shaded}
\begin{Highlighting}[]
\NormalTok{compteur }\OperatorTok{=} \DecValTok{0}

\KeywordTok{def}\NormalTok{ evaluer():}
    \KeywordTok{global}\NormalTok{ compteur}
\NormalTok{    compteur }\OperatorTok{+=} \DecValTok{1}
\end{Highlighting}
\end{Shaded}

\hypertarget{glossaire}{%
\subsection{Glossaire}\label{glossaire}}

dictionnaire dictionary Une collection de paires clé-valeur qui associe
des clés à des valeurs.

clé key Un objet qui apparaît dans un dictionnaire associé à une valeur.

paire clé-valeur key-value pair La représentation d\textquotesingle un
élément dans un dictionnaire.

table de hachage hash table L\textquotesingle algorithme sous-jacent
utilisé pour implémenter les dictionnaires en Python.

hachable hashable Un type qui possède une fonction de hachage. Les types
immuables sont hachables.

mémo memo Un dictionnaire utilisé pour stocker les résultats de calculs
déjà effectués afin d\textquotesingle éviter le double travail.

variable globale global variable Une variable définie en dehors de
toutes les fonctions, accessible partout.

\hypertarget{exercices}{%
\subsection{Exercices}\label{exercices}}

\textbf{Exercice 1}

Écrivez une fonction qui lit les mots d\textquotesingle un fichier et
les stocke comme clés dans un dictionnaire. Utilisez ensuite
l\textquotesingle opérateur \texttt{in} pour vérifier rapidement et
efficacement si un mot est présent dans le dictionnaire.

\textbf{Exercice 2}

Utilisez la méthode \texttt{setdefault} pour écrire une version plus
concise de la fonction \texttt{histogramme} présentée dans ce chapitre.

\textbf{Exercice 3}

Modifiez la fonction \texttt{inverser\_dictionnaire} pour
qu\textquotesingle elle utilise \texttt{setdefault} au lieu
d\textquotesingle une instruction conditionnelle.

\textbf{Exercice 4}

Si vous avez une fonction qui prend une liste et utilise la méthode
\texttt{in} pour trouver les doublons, la complexité temporelle augmente
rapidement. Écrivez une version plus rapide qui utilise un dictionnaire
pour éliminer les doublons d\textquotesingle une liste de manière
efficace.
